\documentclass[../../main.tex]{subfiles}
\begin{document}

引入光滑结构的主要原因，是使我们能够在流形上定义光滑函数以及流形之间的光滑映射。本章我们就来完成这项工作。

我们首先定义光滑的实值函数和向量值函数，然后把它推广到流形之间的光滑映射。接着我们暂时把注意力集中在微分同胚这一特殊情形上，它们是具有光滑逆的双射光滑映射。如果两个光滑流形之间存在一个微分同胚，我们就说它们微分同胚。光滑流形理论的主要研究对象，是那些在微分同胚下不变的性质。

在本章末尾，我们引入一个把局部定义的光滑对象"混合在一起"的强有力工具，称为单位分解。它在整个光滑流形理论中被用来从局部对象构造整体光滑对象。

\section{光滑函数与光滑映射}\label{section:ch02-光滑函数与光滑映射}

虽然在术语上 function（函数）和 map（映射）严格地说同义，但在研究光滑流形时，对它们作一点区分常常是方便的。在全书中，我们一般保留 function（函数）一词指余域为 $\mathbb{R}$（实值函数）或对某个 $k>1$ 为 $\mathbb{R}^{k}$（向量值函数）的映射。而 map 或 mapping 两词都可以指任意类型的映射，例如任意流形之间的映射。

\subsection{流形上的光滑函数}\label{subsection:ch02-流形上的光滑函数}

设 $M$ 是一个光滑 $n$ 流形，$k$ 是非负整数，$f:M\to\mathbb{R}^{k}$ 是任意函数。我们称 $f$ 是一个光滑函数，如果对每个 $p\in M$，存在 $M$ 的一个光滑图 $(U,\varphi)$，其定义域包含 $p$，且使复合函数 $f\circ\varphi^{-1}$ 在开子集 $\widetilde{U}=\varphi(U)\subset\mathbb{R}^{n}$ 上光滑（图~\cref{figure:ch02-2.1}）。若 $M$ 是一个光滑带边流形，定义完全相同，只是现在 $\varphi(U)$ 是 $\mathbb{R}^{n}$ 或 $H^{n}$ 的一个开子集；在后一情形，我们把 $f\circ\varphi^{-1}$ 的光滑性理解为：$\varphi(U)$ 的每一点都有一个（在 $\mathbb{R}^{n}$ 中的）邻域，使 $f\circ\varphi^{-1}$ 在其上延拓为通常意义下的光滑函数。

\begin{figure}[ht]
\centering
\includegraphics[width=0.8\textwidth]{fig-02-1.png}
\caption{光滑函数的定义}\label{figure:ch02-2.1}
\end{figure}

最重要的特殊情形是光滑实值函数 $f:M\to\mathbb{R}$；所有这类函数所成集合记作 $C^{\infty}(M)$。因为光滑函数的和与常数倍仍是光滑的，$C^{\infty}(M)$ 是 $\mathbb{R}$ 上的一个向量空间。

\begin{exercise}\label{exercise:ch02-2.1}
设 $M$ 是一个带边或无边的光滑流形。证明逐点乘法使 $C^{\infty}(M)$ 成为一个交换环，以及 $\mathbb{R}$ 上的一个交换结合代数。（代数的定义见附录~B，第~624 页。）
\end{exercise}

\begin{exercise}\label{exercise:ch02-2.2}
设 $U$ 是 $\mathbb{R}^{n}$ 带其标准光滑流形结构的一个开子流形。证明：函数 $f:U\to\mathbb{R}^{k}$ 在刚才定义的意义下光滑，当且仅当它在通常微积分意义下光滑。对 $H^{n}$ 中一个开的带边子流形也作同样证明（见习题~\cref{exercise:ch01-1.44}）。
\end{exercise}

\begin{exercise}\label{exercise:ch02-2.3}
设 $M$ 是一个带边或无边的光滑流形，并假设 $f:M\to\mathbb{R}^{k}$ 是一个光滑函数。证明对 $M$ 的每个光滑图 $(U,\varphi)$，复合 $f\circ\varphi^{-1}:\varphi(U)\to\mathbb{R}^{k}$ 都光滑。
\end{exercise}

给定函数 $f:M\to\mathbb{R}^{k}$ 和 $M$ 的一个图 $(U,\varphi)$，由 $\widehat{f}(x)=f\circ\varphi^{-1}(x)$ 定义的函数 $\widehat{f}:\varphi(U)\to\mathbb{R}^{k}$ 称为 $f$ 的坐标表示。由定义，$f$ 光滑当且仅当它在每一点附近某个光滑图中的坐标表示光滑。由前面的练习，光滑函数在每个光滑图中都有光滑的坐标表示。

例如，考虑定义在平面上的实值函数 $f(x,y)=x^{2}+y^{2}$。在比如集合 $U=\{(x,y):x>0\}$ 上的极坐标下，它有坐标表示 $\widehat{f}(r,\theta)=r^{2}$。遵照我们用局部坐标把流形的一个开子集等同于欧几里得空间开子集的习惯做法，在不致混淆时我们往往甚至不区分 $\widehat{f}$ 与 $f$ 本身，而代之以说类似"$f$ 在 $U$ 上光滑，因为它的坐标表示 $f(r,\theta)=r^{2}$ 光滑"这样的话。
\subsection{流形之间的光滑映射}\label{subsection:ch02-流形之间的光滑映射}

光滑函数的定义容易推广到流形之间的映射。设 $M,N$ 是光滑流形，$F:M\to N$ 是任意映射。我们称 $F$ 是一个光滑映射，如果对每个 $p\in M$，存在包含 $p$ 的光滑图 $(U,\varphi)$ 和包含 $F(p)$ 的光滑图 $(V,\psi)$，使 $F(U)\subset V$，且复合映射 $\psi\circ F\circ\varphi^{-1}$ 从 $\varphi(U)$ 到 $\psi(V)$ 光滑（图~\cref{figure:ch02-2.2}）。若 $M$ 和 $N$ 都是光滑带边流形，$F$ 的光滑性以完全相同方式定义，照例理解：定义域为 $H^{n}$ 子集的映射在每点附近有到光滑映射的延拓时为光滑；余域为 $H^{n}$ 子集的映射作为映到 $\mathbb{R}^{n}$ 的映射光滑时为光滑。注意，前面实值或向量值函数光滑性的定义，可以看作这一定义取 $N=V=\mathbb{R}^{k}$、$\psi=\operatorname{Id}_{\mathbb{R}^{k}}$ 的特殊情形。

\begin{figure}[ht]
\centering
\includegraphics[width=0.8\textwidth]{fig-02-2.png}
\caption{光滑映射的定义}\label{figure:ch02-2.2}
\end{figure}

关于我们对光滑映射的定义，第一个重要观察是：正如人们所期望的，光滑性蕴涵连续性。

\begin{proposition}\label{proposition:ch02-光滑映射连续}
每个光滑映射都是连续的。
\end{proposition}

\begin{proof}
设 $M$ 和 $N$ 是带边或无边的光滑流形，$F:M\to N$ 光滑。给定 $p\in M$，$F$ 的光滑性意味着存在包含 $p$ 的光滑图 $(U,\varphi)$ 和包含 $F(p)$ 的光滑图 $(V,\psi)$，使 $F(U)\subset V$，且 $\psi\circ F\circ\varphi^{-1}:\varphi(U)\to\psi(V)$ 光滑，从而连续。由于 $\varphi:U\to\varphi(U)$ 和 $\psi:V\to\psi(V)$ 是同胚，这又蕴涵
\begin{align*}
F|_{U}=\psi^{-1}\circ\bigl(\psi\circ F\circ\varphi^{-1}\bigr)\circ\varphi:U\to V
\end{align*}
是连续映射的复合。由于 $F$ 在每点的一个邻域内连续，故它在 $M$ 上连续。
\end{proof}

要直接从定义证明映射 $F:M\to N$ 光滑，部分地要求对每个 $p\in M$，我们证明存在包含 $p$ 的坐标域 $U$ 和包含 $F(p)$ 的坐标域 $V$ 使 $F(U)\subset V$。这一要求被纳入定义，正是为了使光滑性自动蕴涵连续性。（习题~2-1 说明了如果略去这一要求会出什么问题。）

还有其他刻画流形之间映射光滑性的方式，达到同样目的。下面是其中两种。

\begin{proposition}[光滑性的等价刻画]\label{proposition:ch02-光滑性等价刻画}
设 $M,N$ 是带边或无边的光滑流形，$F:M\to N$ 是一个映射。则 $F$ 光滑，当且仅当满足以下条件之一：
\begin{enumerate}[(a)]
\item 对每个 $p\in M$，存在包含 $p$ 的光滑图 $(U,\varphi)$ 和包含 $F(p)$ 的光滑图 $(V,\psi)$，使 $U\cap F^{-1}(V)$ 在 $M$ 中开，且复合映射 $\psi\circ F\circ\varphi^{-1}$ 从 $\varphi(U\cap F^{-1}(V))$ 到 $\psi(V)$ 光滑。
\item $F$ 连续，且存在 $M$ 和 $N$ 的光滑图册 $\{(U_{\alpha},\varphi_{\alpha})\}$ 和 $\{(V_{\beta},\psi_{\beta})\}$，使对每个 $\alpha$ 和 $\beta$，映射 $\psi_{\beta}\circ F\circ\varphi_{\alpha}^{-1}$ 从 $\varphi_{\alpha}(U_{\alpha}\cap F^{-1}(V_{\beta}))$ 到 $\psi_{\beta}(V_{\beta})$ 光滑。
\end{enumerate}
\end{proposition}

\begin{proposition}[光滑性是局部的]\label{proposition:ch02-光滑性局部}
设 $M,N$ 是带边或无边的光滑流形，$F:M\to N$ 是一个映射。
\begin{enumerate}[(a)]
\item 若 $M$ 的每个点 $p$ 都有一个邻域 $U$ 使限制 $F|_{U}$ 光滑，则 $F$ 光滑。
\item 反之，若 $F$ 光滑，则它在每个开子集上的限制都光滑。
\end{enumerate}
\end{proposition}

\begin{exercise}\label{exercise:ch02-2.7}
证明上面两个命题。
\end{exercise}

下一个推论本质上只是前一命题的重述，但它给出了构造光滑映射的一种极有用的方式。

\begin{corollary}[光滑映射的粘合引理]\label{corollary:ch02-光滑映射粘合引理}
设 $M,N$ 是带边或无边的光滑流形，$\{U_{\alpha}\}_{\alpha\in A}$ 是 $M$ 的一个开覆盖。假设对每个 $\alpha\in A$，给定一个光滑映射 $F_{\alpha}:U_{\alpha}\to N$，使这些映射在重叠处一致：对所有 $\alpha$ 和 $\beta$ 有 $F_{\alpha}|_{U_{\alpha}\cap U_{\beta}}=F_{\beta}|_{U_{\alpha}\cap U_{\beta}}$。则存在唯一的光滑映射 $F:M\to N$，使对每个 $\alpha\in A$ 有 $F|_{U_{\alpha}}=F_{\alpha}$。
\end{corollary}

若 $F:M\to N$ 是一个光滑映射，$(U,\varphi)$ 和 $(V,\psi)$ 分别是 $M$ 和 $N$ 的任意光滑图，我们称 $\widehat{F}=\psi\circ F\circ\varphi^{-1}$ 为 $F$ 关于所给坐标的坐标表示。它把集合 $\varphi(U\cap F^{-1}(V))$ 映到 $\psi(V)$。

\begin{exercise}\label{exercise:ch02-2.9}
假设 $F:M\to N$ 是带边或无边光滑流形之间的一个光滑映射。证明 $F$ 关于 $M$ 和 $N$ 的每一对光滑图的坐标表示都光滑。
\end{exercise}

与实值或向量值函数一样，一旦我们在定义域和余域中都选定了具体的局部坐标，就常常可以忽略 $F$ 与 $\widehat{F}$ 之间的区别。

接下来我们考察几类自动光滑的简单映射。

\begin{proposition}\label{proposition:ch02-自动光滑映射}
设 $M,N,P$ 是带边或无边的光滑流形。
\begin{enumerate}[(a)]
\item 每个常值映射 $c:M\to N$ 都光滑。
\item $M$ 的恒等映射光滑。
\item 若 $U\subset M$ 是一个带边或无边的开子流形，则包含映射 $U\hookrightarrow M$ 光滑。
\item 若 $F:M\to N$ 和 $G:N\to P$ 光滑，则 $G\circ F:M\to P$ 也光滑。
\end{enumerate}
\end{proposition}

\begin{proof}
我们证明 (d)，其余留作练习。设 $F:M\to N$ 和 $G:N\to P$ 是光滑映射，$p\in M$。由 $G$ 光滑性的定义，存在包含 $F(p)$ 的光滑图 $(V,\phi)$ 和包含 $G(F(p))$ 的光滑图 $(W,\theta)$，使 $G(V)\subset W$，且 $\theta\circ G\circ\phi^{-1}:\phi(V)\to\theta(W)$ 光滑。由于 $F$ 连续，$F^{-1}(V)$ 是 $M$ 中 $p$ 的一个邻域；于是存在 $M$ 的一个光滑图 $(U,\varphi)$ 使 $p\in U\subset F^{-1}(V)$（图~\cref{figure:ch02-2.3}）。由习题~\cref{exercise:ch02-2.9}，$\phi\circ F\circ\varphi^{-1}$ 从 $\varphi(U)$ 到 $\phi(V)$ 光滑。于是有 $G\circ F(U)\subset G(V)\subset W$，且
\begin{align*}
\theta\circ(G\circ F)\circ\varphi^{-1}=\bigl(\theta\circ G\circ\phi^{-1}\bigr)\circ\bigl(\phi\circ F\circ\varphi^{-1}\bigr):\varphi(U)\to\theta(W)
\end{align*}
光滑，因为它是欧几里得空间子集之间光滑映射的复合。
\end{proof}

\begin{figure}[ht]
\centering
\includegraphics[width=0.8\textwidth]{fig-02-3.png}
\caption{光滑映射的复合是光滑的}\label{figure:ch02-2.3}
\end{figure}

\begin{exercise}\label{exercise:ch02-2.11}
证明上面命题的 (a)--(c) 部分。
\end{exercise}

\begin{proposition}\label{proposition:ch02-乘积映射分量}
假设 $M_{1},\dots,M_{k}$ 和 $N$ 是带边或无边的光滑流形，使得 $M_{1},\dots,M_{k}$ 中至多一个有非空边界。对每个 $i$，令 $\pi_{i}:M_{1}\times\cdots\times M_{k}\to M_{i}$ 表示到第 $i$ 个因子的投影。映射 $F:N\to M_{1}\times\cdots\times M_{k}$ 光滑，当且仅当每个分量映射 $F_{i}=\pi_{i}\circ F:N\to M_{i}$ 都光滑。
\end{proposition}

\begin{proof}
习题~2-2。
\end{proof}
尽管本书的大部分精力用于研究光滑流形和光滑映射，但我们偶尔也需要与拓扑流形和连续映射打交道。为一致起见，我们采用如下约定：不加进一步限定时，"函数"和"映射"两词纯粹在集合论意义下理解，不带连续性或光滑性的任何假设。然而，我们研究的大多数其他对象，默认将被理解为带有某种最小拓扑结构。除非另行说明，"流形"或"带边流形"总理解为拓扑的，而"坐标图"在拓扑意义下理解，即从流形的一个开子集到 $\mathbb{R}^{n}$ 或 $H^{n}$ 的一个开子集的同胚。如果我们希望把注意力限制在光滑流形或光滑坐标图上，我们会明说。类似地，我们对许多其他特定类型的几何对象及其间映射的默认假设，至多是连续性；除非明确指定，不假定光滑性。唯一的例外是少数一些概念，它们的定义本身就要求光滑性。

这一约定要求某种纪律，因为每当需要光滑性假设时我们都得记得陈述它；但其好处是它（在很大程度上）使我们免于去记哪几类映射假定是光滑的、哪几类不是。

另一方面，因为光滑映射的定义要求定义域和余域都有光滑结构，如果我们说"$F:M\to N$ 是一个光滑映射"而未具体说明 $M$ 和 $N$ 是什么，应始终理解它们是带边或无边的光滑流形。

我们现在已有足够信息来给出若干有趣的光滑映射例子。尽管定义表面上复杂，证明某个具体映射光滑通常并不困难。基本上只有三种常见做法：
\begin{itemize}
\item 用光滑局部坐标写出该映射，并认出其分量函数是光滑初等函数的复合；
\item 把该映射表示为已知光滑映射的复合；
\item 使用适用于所考虑具体情形的某个专门定理。
\end{itemize}

\begin{example}[光滑映射]\label{example:ch02-光滑映射}
\begin{enumerate}[(a)]
\item 从一个 $0$ 维流形映到带边或无边光滑流形的任意映射都自动光滑，因为每个坐标表示都是常值的。
\item 若圆 $S^{1}$ 被赋予其标准光滑结构，则由 $\varepsilon(t)=e^{2\pi\mathrm{i}t}$ 定义的映射 $\varepsilon:\mathbb{R}\to S^{1}$ 光滑，因为关于 $S^{1}$ 的任意角坐标 $\theta$（见习题~\cref{exercise:ch01-1-8}），它有形如
\begin{align*}
\widehat{\varepsilon}(t)=2\pi t+c
\end{align*}
的坐标表示，其中 $c$ 是某个常数，你可以自行验证。
\item 由
\begin{align*}
\varepsilon_{n}\bigl(x^{1},\dots,x^{n}\bigr)=\bigl(e^{2\pi\mathrm{i}x^{1}},\dots,e^{2\pi\mathrm{i}x^{n}}\bigr)
\end{align*}
定义的映射 $\varepsilon_{n}:\mathbb{R}^{n}\to T^{n}$ 由命题~\cref{proposition:ch02-乘积映射分量} 是光滑的。
\item 现在考虑带其标准光滑结构的 $n$ 球面 $S^{n}$。包含映射 $\iota:S^{n}\hookrightarrow\mathbb{R}^{n+1}$ 当然连续，因为它是一个拓扑子空间的包含映射。它是一个光滑映射，因为关于例~\cref{example:ch01-球面光滑结构} 的任意图像坐标，它的坐标表示为
\begin{align*}
\widehat{\iota}\bigl(u^{1},\dots,u^{n}\bigr)=\iota\circ(\varphi_{i}^{\pm})^{-1}\bigl(u^{1},\dots,u^{n}\bigr)
=\bigl(u^{1},\dots,u^{i-1},\pm\sqrt{1-|u|^{2}},u^{i},\dots,u^{n}\bigr),
\end{align*}
这在其定义域（$|u|^{2}<1$ 的集合）上是光滑的。
\item 用来定义 $\mathbb{RP}^{n}$ 的商映射 $\pi:\mathbb{R}^{n+1}\setminus\{0\}\to\mathbb{RP}^{n}$ 是光滑的，因为它用例~\cref{example:ch01-射影空间光滑结构} 构造的 $\mathbb{RP}^{n}$ 坐标与 $\mathbb{R}^{n+1}\setminus\{0\}$ 上标准坐标表示的坐标表示为
\begin{align*}
\widehat{\pi}\bigl(x^{1},\dots,x^{n+1}\bigr)=\varphi_{i}\bigl([x^{1},\dots,x^{n+1}]\bigr)
=\left(\frac{x^{1}}{x^{i}},\dots,\frac{x^{i-1}}{x^{i}},\frac{x^{i+1}}{x^{i}},\dots,\frac{x^{n+1}}{x^{i}}\right).
\end{align*}
\item 把 $\pi:\mathbb{R}^{n+1}\setminus\{0\}\to\mathbb{RP}^{n}$ 限制到 $S^{n}\subset\mathbb{R}^{n+1}\setminus\{0\}$ 来定义 $q:S^{n}\to\mathbb{RP}^{n}$。它是一个光滑映射，因为它是前两个例子中映射的复合 $q=\pi\circ\iota$。
\item 若 $M_{1},\dots,M_{k}$ 是光滑流形，则每个投影映射 $\pi_{i}:M_{1}\times\cdots\times M_{k}\to M_{i}$ 都光滑，因为关于例~\cref{example:ch01-乘积流形} 的任意乘积坐标，它的坐标表示恰为一个坐标投影。
\end{enumerate}
\end{example}

\subsection{微分同胚}\label{subsection:ch02-微分同胚}

若 $M,N$ 是带边或无边的光滑流形，从 $M$ 到 $N$ 的一个微分同胚是一个光滑双射 $F:M\to N$，且具有光滑逆。如果两个光滑流形之间存在一个微分同胚，我们就说它们微分同胚。有时记作 $M\cong N$。

\begin{example}[微分同胚]\label{example:ch02-微分同胚}
\begin{enumerate}[(a)]
\item 考虑由
\begin{align}\label{align:ch02-Bn与Rn微分同胚}
F(x)=\frac{x}{\sqrt{1-|x|^{2}}},\qquad G(y)=\frac{y}{\sqrt{1+|y|^{2}}}
\end{align}
给出的映射 $F:B^{n}\to\mathbb{R}^{n}$ 和 $G:\mathbb{R}^{n}\to B^{n}$。这些映射光滑，并且直接计算可知它们互为逆映射。因此它们都是微分同胚，从而 $B^{n}$ 微分同胚于 $\mathbb{R}^{n}$。
\item 若 $M$ 是任意光滑流形，$(U,\varphi)$ 是 $M$ 上的一个光滑坐标图，则 $\varphi:U\to\varphi(U)\subset\mathbb{R}^{n}$ 是一个微分同胚。（事实上，它以恒等映射为坐标表示。）
\end{enumerate}
\end{example}

\begin{proposition}[微分同胚的性质]\label{proposition:ch02-微分同胚性质}
\begin{enumerate}[(a)]
\item 微分同胚的每个复合都是微分同胚。
\item 光滑流形之间微分同胚的每个有限乘积都是微分同胚。
\item 每个微分同胚都是同胚且是开映射。
\item 一个微分同胚在一个带边或无边开子流形上的限制，是到其像的微分同胚。
\item "微分同胚"是所有带边或无边光滑流形所成类上的一个等价关系。
\end{enumerate}
\end{proposition}

\begin{exercise}\label{exercise:ch02-2.16}
证明上面的命题。
\end{exercise}

下面的定理是维数不变性的一个弱版本，对许多目的已经足够。

\begin{theorem}[维数的微分同胚不变性]\label{theorem:ch02-维数微分同胚不变性}
一个维数为 $m$ 的非空光滑流形不可能微分同胚于一个 $n$ 维光滑流形，除非 $m=n$。
\end{theorem}

\begin{proof}
设 $M$ 是一个非空光滑 $m$ 流形，$N$ 是一个非空光滑 $n$ 流形，$F:M\to N$ 是一个微分同胚。任取一点 $p\in M$，令 $(U,\varphi)$ 和 $(V,\psi)$ 分别是包含 $p$ 和 $F(p)$ 的光滑坐标图。则（其限制）$\widehat{F}=\psi\circ F\circ\varphi^{-1}$ 是从 $\mathbb{R}^{m}$ 的一个开子集到 $\mathbb{R}^{n}$ 的一个开子集的微分同胚，因此由命题~C.4 得 $m=n$。
\end{proof}

对边界也有类似的不变性命题。

\begin{theorem}[边界的微分同胚不变性]\label{theorem:ch02-边界微分同胚不变性}
设 $M,N$ 是光滑带边流形，$F:M\to N$ 是一个微分同胚。则 $F(\partial M)=\partial N$，且 $F$ 限制为一个从 $\operatorname{Int}M$ 到 $\operatorname{Int}N$ 的微分同胚。
\end{theorem}

\begin{exercise}\label{exercise:ch02-2.19}
利用定理~\cref{theorem:ch01-边界光滑不变性} 证明上面的定理。
\end{exercise}

正如两个拓扑空间如果同胚就被认为是"相同的"一样，两个带边或无边的光滑流形如果微分同胚，本质上是不可区分的。光滑流形理论的核心关切，是研究那些被微分同胚保持的光滑流形性质。定理~\cref{theorem:ch02-维数微分同胚不变性} 表明维数就是这样一种性质。

人们自然会问，给定拓扑流形上的光滑结构是否唯一。这个问题的这种直截了当版本容易回答：我们在例~\cref{example:ch01-0维流形} 中已观察到每个 $0$ 维流形有唯一的光滑结构，但正如习题~1-6 所示，每个正维流形一旦容许一个光滑结构，就容许许多不同的光滑结构。

一个更微妙也更有趣的问题是：给定拓扑流形是否容许彼此不微分同胚的光滑结构。例如，令 $\widetilde{\mathbb{R}}$ 表示拓扑流形 $\mathbb{R}$，但赋予例~\cref{example:ch01-另一光滑结构} 中描述的光滑结构（由整体图 $\psi(x)=x^{3}$ 定义）。结果表明 $\widetilde{\mathbb{R}}$ 微分同胚于带其标准光滑结构的 $\mathbb{R}$。定义一个映射 $F:\mathbb{R}\to\widetilde{\mathbb{R}}$ 为 $F(x)=x^{1/3}$。这个映射的坐标表示为 $\widehat{F}(t)=\psi\circ F\circ\operatorname{Id}_{\mathbb{R}}^{-1}(t)=t$，显然光滑。此外，其逆的坐标表示为
\begin{align*}
\widehat{F^{-1}}(y)=\operatorname{Id}_{\mathbb{R}}\circ F^{-1}\circ\psi^{-1}(y)=y,
\end{align*}
也光滑，所以 $F$ 是一个微分同胚。（这是一个重要情形：必须保持映射与其坐标表示之间的区别！）

事实上，正如你后面会看到的，$\mathbb{R}$ 上在微分同胚意义下只有一个光滑结构（见习题~15-13）。更精确地说，若 $A_{1}$ 和 $A_{2}$ 是 $\mathbb{R}$ 上任意两个光滑结构，则存在一个微分同胚 $F:(\mathbb{R},A_{1})\to(\mathbb{R},A_{2})$。事实上，由 James Munkres [Mun60] 和 Edwin Moise [Moi77] 的工作可得，每个维数小于或等于 $3$ 的拓扑流形都有一个在微分同胚意义下唯一的光滑结构。高维中的类似问题结果相当深刻，且很大程度上仍未解决。即使对于欧几里得空间，光滑结构唯一性的问题直到二十世纪后期才完全解决。答案令人惊讶：只要 $n\neq 4$，$\mathbb{R}^{n}$ 有唯一的光滑结构（在微分同胚意义下）；但 $\mathbb{R}^{4}$ 有不可数个不同的光滑结构，其中任意两个都不彼此微分同胚！$\mathbb{R}^{4}$ 上非标准光滑结构（称为"怪 $\mathbb{R}^{4}$"）的存在性，由 Simon Donaldson 和 Michael Freedman 于 1984 年作为他们关于紧 $4$ 流形几何与拓扑工作的结果首先证明；结果描述于 [DK90] 和 [FQ90]。

对紧流形，情况更加迷人。1956 年，John Milnor [Mil56] 证明 $S^{7}$ 上存在不与标准结构微分同胚的光滑结构。后来，他和 Michel Kervaire [KM63]（利用 Steve Smale 的一个深刻定理 [Sma62]）证明这类结构恰有 $15$ 个微分同胚类（如果你限制为保持一个称为定向的性质的微分同胚，则有 $28$ 类，定向将在第~15 章讨论）。另一方面，在所有大于 $3$ 的维数中，存在根本没有任何光滑结构的紧拓扑流形。确定拓扑 $4$ 流形上光滑结构数目（若有的话）的问题，是当前研究的一个活跃课题。

\end{document}